\chapter*{Anexo E: Registro de limitaciones}
\addcontentsline{toc}{chapter}{Anexo E: Registro de limitaciones}
\label{cap:anexo_limitaciones}

\setcounter{section}{0}
\renewcommand{\thesection}{E.\arabic{section}}

Este anexo recoge las limitaciones que el desarrollo de la plataforma dejó
documentadas, de las que la Sección~\ref{sec:conclusiones_limitaciones} del
Capítulo~\ref{cap:conclusiones} consolida solo las de mayor alcance. El
registro mantenido en el repositorio, en
\path{docs/pipeline/known_limitations.md}, es la fuente autoritativa y
describe de cada limitación su impacto, su tratamiento actual y su estado.
Este anexo las agrupa en cuatro familias y indica, cuando existe, cómo se
resolvió o se acotó cada una.

\section{Infraestructura y cadena de suministro}
\label{sec:anexo_limitaciones_infra}

\begin{itemize}
  \item \textbf{Imágenes de Bitnami congeladas.} Las imágenes de MinIO, de
  las dos instancias de PostgreSQL y de los subgráficos de Superset se fijan
  al registro congelado de Bitnami y no reciben parches de seguridad. Se
  acepta porque el clúster nunca se expone a Internet.
  \item \textbf{Chart de Superset obsoleto.} Sus mantenedores lo han marcado
  como obsoleto y la vía oficial es ahora un operador de Kubernetes. Se
  acepta, y el operador queda como trabajo futuro.
  \item \textbf{Superset local y sin endurecer.} Tiene una réplica de cada
  componente, sin TLS ni inicio de sesión único, y sus metadatos no tienen
  copia de seguridad. Se acepta para un clúster local.
  \item \textbf{Montajes del nodo e imagen importada a mano.} Los flujos de
  ingesta y transformación montan código y datos desde el repositorio y
  contienen la ruta absoluta de la máquina de desarrollo. Es una decisión
  deliberada que evita reconstruir imágenes al editar código, y limita los
  flujos a un único nodo.
  \item \textbf{Python 3.10 en el trabajo de silver.} El trabajo corre con el
  intérprete de la imagen de Spark y no con el 3.14 del repositorio, de modo
  que el código compartido debe seguir siendo compatible con 3.10.
  \item \textbf{Registro de cliente en la API de ORCID bloqueado.} Un
  proyecto individual no puede registrar una aplicación cliente. La
  plataforma usa el nivel anónimo, con un límite diario menor.
  \item \textbf{Propiedades del controlador de Spark sin efecto.} En modo
  cliente, la configuración de pod del controlador no se aplica y
  \texttt{spark-submit} no avisa. La configuración se fija en el propio pod
  que lanza el trabajo.
  \item \textbf{Borrado por lotes como permiso aparte.} Un rol que concede
  borrado pero no borrado de colección deja al controlador de Spark sin poder
  limpiar sus pods. El rol del repositorio concede ambos.
  \item \textbf{Servidor de API de Airflow cerrado por su sonda de vida.}
  Bajo carga, la sonda mataba el servidor y las tareas que arrancaban en ese
  intervalo fallaban. Resuelta con recursos y sondas explícitos.
  \item \textbf{Reinicio bloqueado del planificador de Airflow.} Una fila de
  tarea huérfana, dejada por un planificador caído, hace fallar todos los
  reinicios siguientes por un error aún abierto en Airflow. Resuelta
  operativamente marcando como fallidas esa tarea y su ejecución.
  \item \textbf{Pods ejecutores en estado desconocido.} Un planificador caído
  puede dejar pods sin limpiar. Se documenta junto con la limitación anterior.
  \item \textbf{Ejecución perdida al activar un flujo.} Activar un flujo con
  planificación diaria crea de inmediato la ejecución más reciente que se
  perdió. Es inocuo porque los flujos son idempotentes por identificador de
  ejecución.
  \item \textbf{Primer arranque de PostgreSQL interrumpido.} Si se interrumpe
  durante su inicialización, nunca crea el rol de gold y no se recupera solo.
  Resuelta operativamente borrando el pod y su volumen.
  \item \textbf{Memoria insuficiente de PostgreSQL.} El perfil mínimo del
  \textit{chart} mataba el contenedor al preparar las tablas de mayor
  volumen. Resuelta con recursos explícitos.
  \item \textbf{Pruebas de Spark en contenedores.} Con demasiados
  contenedores a la vez, la batería fallaba por tiempo agotado. Resuelta
  limitando a cuatro los contenedores simultáneos.
\end{itemize}

\section{Alcance de los datos y de la ingesta}
\label{sec:anexo_limitaciones_datos}

\begin{itemize}
  \item \textbf{Subconjunto de ORCID puntual.} Es una instantánea filtrada
  por país de afiliación, con supuestos de serialización.
  \item \textbf{Muestra acotada del subconjunto.} Por defecto se aterrizan
  los primeros 20\,000 registros en orden de cubo, que ocupan 2,7 GB, y no
  el subconjunto completo. El parámetro del flujo permite ampliarlo.
  \item \textbf{Datos sintéticos mixtos.} Los CVN sintéticos combinan datos
  públicos reales de ORCID con datos personales inventados.
  \item \textbf{Cobertura estrecha del generador.} Solo produce cuatro tipos
  de entrada, que son identidad, experiencia profesional, formación y
  publicaciones.
  \item \textbf{Extracción parcial en silver.} Solo se extraen esos cuatro
  tipos. Las demás secciones se validan sin extraerse, solo se traduce el
  código numérico de país de España y las obras de ORCID no traen autores.
  \item \textbf{Esquema JSON permisivo.} No impone la forma de cada entrada y
  difiere del importador XML en el tratamiento de las fechas. Se hereda del
  trabajo anterior.
  \item \textbf{Lectura conjunta de bronze.} Leerlo como un solo conjunto
  mezcla cargas útiles de texto y de objeto. Cada fuente se lee por su propia
  ruta.
  \item \textbf{Cifras de carrera como cotas inferiores.} El 18\% de las
  afiliaciones no tiene año de inicio y el 33\% no tiene año de fin, y una
  estancia nunca se extiende hasta hoy.
\end{itemize}

\section{Resolución de identidad e indicadores}
\label{sec:anexo_limitaciones_resolucion}

\begin{itemize}
  \item \textbf{Resolución conservadora.} La regla de nombre y afiliación
  alcanza una precisión del 95,8\% y una cobertura del 74,2\%, por debajo del
  objetivo del 99\% de precisión. Una persona con dos identificadores ORCID
  cuenta como dos entidades.
  \item \textbf{Calidad medida solo contra datos sintéticos.} La verdad de
  referencia es el identificador de la semilla y solo hay cuatro variantes de
  nombre, de modo que las cifras son optimistas para currículos reales.
  \item \textbf{Colaboración como cota inferior.} Silver no tiene entidades
  de coautor, de modo que dos entidades solo se enlazan si declaran el mismo
  DOI. Los 9\,540 pares incluyen 25 que son una misma persona vista dos
  veces.
  \item \textbf{Deduplicación por tipo de clave.} Una obra con DOI en un
  registro y sin él en otro cuenta dos veces. El 1,2\% de las obras no tiene
  año útil y queda fuera del indicador por año.
  \item \textbf{Errores de resolución heredados.} Una fusión falsa mezcla las
  publicaciones de dos personas. Deshacer las 75 fusiones de una sola
  organización cambia las entidades un 0,25\% y los pares de colaboración un
  0,58\%.
  \item \textbf{Panel sobre agregados parciales.} Deja fuera 443 obras
  anteriores a 1980 y 16\,231 de 94\,431 estancias sin año de inicio o fuera
  del intervalo 1970-2026, y desactiva su caché de datos.
\end{itemize}

\section{Escalabilidad y medición}
\label{sec:anexo_limitaciones_escala}

\begin{itemize}
  \item \textbf{Reconstrucción completa.} Silver y gold se recalculan por
  entero en cada ejecución y PostgreSQL conserva solo la última publicación.
  Es adecuado para 20\,000 registros masivos y no para los 301\,763 del
  subconjunto completo.
  \item \textbf{Techo de memoria con pocos ejecutores.} Con el tamaño de
  memoria fijo, un ejecutor, y a cuádruple volumen también dos, agotan la
  memoria del trabajo de silver de forma determinista.
  \item \textbf{Aceleración limitada en un solo nodo.} Más ejecutores no
  reducen de forma fiable el tiempo, porque el nodo no está limitado por el
  procesador.
\end{itemize}
